\documentclass[10pt]{article}
\usepackage[T1]{fontenc}
\usepackage[utf8]{inputenc}
\usepackage[french]{babel}
\usepackage{amsmath,amssymb,mathrsfs,mathpazo}
\usepackage{geometry}
\geometry{a4paper,margin=2cm}
\usepackage{graphicx,array,tabularx,tabularray,colortbl,multirow,multicol,booktabs,enumitem,pifont,ulem,fancybox,diagbox,cancel}
\usepackage{pstricks,pstricks-add,pst-plot,pst-tree,pst-node,pst-text,pst-func,pst-math,pst-bspline,pst-3dplot,pst-eucl,pst-solides3d}
\let\pscalculate\undefined
\usepackage{pst-bezier}
\usepackage{tikz,pgfplots,tkz-tab}
\usetikzlibrary{arrows,arrows.meta,patterns,calc,positioning,trees,decorations.markings,shapes.geometric,shapes.misc,shapes.arrows,intersections,decorations.pathmorphing,decorations.pathreplacing,backgrounds,shadows,babel}
\pgfplotsset{compat=1.18}
\usepgfplotslibrary{fillbetween}
\usepackage[np]{numprint}
\usepackage{listings,makecell,pas-tableur,twoopt}
\let\barre\undefined
\usepackage[pstricks]{tablvar}
\definecolor{darkgreen}{RGB}{0,100,0}
\providecommand{\red}{\color{red}}
\providecommand{\blue}{\color{blue}}
\providecommand{\green}{\color{green}}
\providecommand{\black}{\color{black}}
\providecommand{\white}{\color{white}}
\providecommand{\gray}{\color{gray}}
\providecommand{\cyan}{\color{cyan}}
\providecommand{\magenta}{\color{magenta}}
\providecommand{\yellow}{\color{yellow}}
\providecommand{\violet}{\color{violet}}
\providecommand{\orange}{\color{orange}}
\providecommand{\euro}{\texteuro}
\providecommand{\EUR}{\texteuro}
\providecommand{\ieme}{\textsuperscript{e}}
\providecommand{\degres}{\ensuremath{^\circ}}
\providecommand{\degree}{\ensuremath{^\circ}}
\providecommand{\vv}[1]{\overrightarrow{#1}}
\providecommand{\Coord}[1]{#1}
\providecommand{\decosix}{}
\providecommand{\starredbullet}{\ensuremath{\bullet}}
\pagestyle{empty}
\setlength{\parindent}{0pt}
\usepackage{adjustbox}\DeclareUnicodeCharacter{2236}{:}\DeclareUnicodeCharacter{2190}{\ensuremath{\leftarrow}}
\providecommand{\R}{\mathbb{R}}
\providecommand{\N}{\mathbb{N}}
\providecommand{\Z}{\mathbb{Z}}
\providecommand{\Q}{\mathbb{Q}}
\providecommand{\D}{\mathbb{D}}
\providecommand{\C}{\mathbb{C}}
\providecommand{\e}{{\mathrm{e}}}
\providecommand{\vect}[1]{\overrightarrow{#1}}
\providecommand{\vectt}[1]{\overrightarrow{\mathrm{#1}}}
\providecommand{\vvt}[1]{\overrightarrow{\mathrm{#1}}}
\providecommand{\np}[2][]{{#2~#1}}
\providecommand{\vv}[1]{{\overrightarrow{#1}}}
\renewcommand{\d}{{\mathrm{d}}}
\providecommand{\degres}{{^\circ}}
\providecommand{\degre}{{^\circ}}
\providecommand{\blue}{\color{blue}}
\providecommand{\red}{\color{red}}
\providecommand{\enskip}{\quad}
\providecommand{\llbracket}{\lbrack\!\lbrack}
\providecommand{\rrbracket}{\rbrack\!\rbrack}
\providecommand{\euro}{{\text{€}}}
\providecommand{\no}{{\text{n°}}}
\providecommand{\interval}[2]{{[#1;#2]}}
\providecommand{\intervFO}[2]{{[#1;#2[}}
\providecommand{\dotfill}{\ldots}

\providecommand{\ds}{\displaystyle}
\providecommand{\textding}[1]{#1}
\providecommand{\Oij}{\(\mathrm{O};\vec{\imath},\vec{\jmath}\)}
\providecommand{\Oijk}{\(\mathrm{O};\vec{\imath},\vec{\jmath},\vec{k}\)}
\providecommand{\Ouv}{\(\mathrm{O};\vec{u},\vec{v}\)}
\providecommand{\intervFF}[2]{\ensuremath{[#1;#2]}}
\renewcommand{\euro}{\ifmmode\text{\texteuro}\else\texteuro\fi}
\DeclareUnicodeCharacter{00B0}{\ifmmode^{\circ}\else\textdegree\fi}
\DeclareUnicodeCharacter{20AC}{\ifmmode\text{\texteuro}\else\texteuro\fi}
% Énoncés : Éducation nationale. Transcription et corrigés : APMEP, Denis Vergès et contributeurs. Découpage et classement : CoopMaths, Sébastien Lozano, Éric Elter et contributeurs. Adaptation de la présentation : Maths Explorer.
% CC BY-SA : https://coopmaths.fr/www/about/
% https://coopmaths.fr/alea/static/eam/2026/tex/eam_specifique_2026_04_polynesie_3_cor.tex
\begin{document}
\noindent\textit{APMEP / CoopMaths. CC BY-SA. Présentation adaptée par Maths Explorer.}\par\medskip
\emph{Pour $t \in [0~;~10]$, $f(t) = -25t^2 + 300t + 400$ où $t$ est le nombre de semaines écoulées
depuis le début de l'épidémie et $f(t)$ le nombre hebdomadaire de consultations.}

\subsection*{Partie A}

\begin{center}
	\begin{tikzpicture}[x=1.1cm,y=0.0056cm,>=stealth]
		\draw[xstep=1,ystep=100,gray!75,line width=0.5pt] (0,0) grid (10.2,1420);
		\draw[->,line width=1pt] (0,0) -- (10.6,0);
		\draw[->,line width=1pt] (0,0) -- (0,1480);
		\foreach \x in {0,1,...,10} \draw (\x,12) -- (\x,-12) node[below] {\footnotesize \np{\x}};
		\foreach \y in {100,200,...,1400} \draw (0.08,\y) -- (-0.08,\y) node[left] {\footnotesize \np{\y}};
		\draw[blue,dashed,line width=1pt] (0,1200) -- (10.2,1200);
		\draw[blue,dashed] (4,1200) -- (4,0);
		\draw[blue,dashed] (8,1200) -- (8,0);
		\draw[blue,line width=2.5pt] (4,0) -- (8,0);
		\draw[dashed] (2,0) -- (2,900) -- (0,900);
		\draw[dashed] (6,0) -- (6,1300) -- (0,1300);
		\draw[red,line width=1.2pt] plot[domain=0:10,samples=80] (\x,{-25*\x*\x+300*\x+400});
		\node[red] at (9.4,1060) {$\mathcal{C}_f$};
	\end{tikzpicture}
\end{center}

\begin{enumerate}
	\item Par lecture graphique, $f(2) \approx 900$.

	Deux semaines après le début de l'épidémie, il y a environ 900 consultations dans la
	semaine.

	\item La courbe coupe la droite d'équation $y = \np{1200}$ aux points d'abscisses 4 et 8, et se situe
	au-dessus entre ces deux valeurs.

	Le nombre de consultations dépasse \np{1200} pour $t \in \left]4~;~8\right[$, c'est-à-dire
	entre la 4\up{e} et la 8\up{e} semaine après le début de
	l'épidémie.

	\item Le sommet de la parabole a pour coordonnées $(6~;~\np{1300})$.

	Le nombre de consultations est maximal au bout de 6 semaines ; il atteint alors
	\np{1300} consultations dans la semaine.
\end{enumerate}

\subsection*{Partie B}

\begin{enumerate}
	\item $f(0) = -25\times 0^2 + 300\times 0 + 400 = 400$.

	Au début de l'épidémie, on compte 400 consultations dans la semaine.

	\item Pour tout réel $t$ de $[0~;~10]$,
$f'(t) = -25\times 2t + 300\times 1 + 0 = -50t + 300$.

	\item $f'(t) = 0 \iff -50t + 300 = 0 \iff -50t = -300 \iff t = \dfrac{-300}{-50} = 6$.
	
$S = \{6\}$.

	\item $f'(t) = -50t + 300$ est une expression affine de coefficient directeur $-50 < 0$ : elle est
	positive avant 6 et négative après.

	$f(0) = 400$, $f(6) = \np{1300}$ et
	$f(10) = -\np{2500} + \np{3000} + 400 = 900$.

	\begin{center}
		\begin{tikzpicture}
			\tkzTabInit[lgt=3.4,espcl=2.8]{$t$/1,$f'(t)$/1,$f$/2}{$0$,$6$,$10$}
			\tkzTabLine{,+,z,-,}
			\tkzTabVar{-/$400$,+/$\np{1300}$,-/$900$}
		\end{tikzpicture}
	\end{center}

	La fonction $f$ est croissante sur $[0~;~6]$ et décroissante sur $[6~;~10]$. Elle admet un
	maximum égal à \np{1300} en $t = 6$.
\end{enumerate}

\end{document}
